% E1.3 -- lema de aproximação em Besov pelo sieve de wavelets periodizadas.
% Transposição do Passo 1 do WALL teórico (ms_theo_1.tex, Lemma "Besov
% approximation error"; prova em supp_theo_1.tex, S.1). Conferência numérica:
% check/02-aproximacao-besov.R (imprime OK). Notação: docs/notacao.md (E1.1).
\documentclass[11pt,a4paper]{article}
\usepackage[T1]{fontenc}
\usepackage[utf8]{inputenc}
\usepackage[brazil]{babel}
\usepackage[margin=2.6cm]{geometry}
\usepackage{enumitem}
\usepackage{hyperref}
% Macros do WAFC. Implementa docs/notacao.md (esboço; congela em E1.1).
% Único lugar onde uma macro é definida; os derivations/*.tex fazem % Macros do WAFC. Implementa docs/notacao.md (esboço; congela em E1.1).
% Único lugar onde uma macro é definida; os derivations/*.tex fazem % Macros do WAFC. Implementa docs/notacao.md (esboço; congela em E1.1).
% Único lugar onde uma macro é definida; os derivations/*.tex fazem \input{macros}.
\usepackage{amsmath,amssymb,amsthm}
\newcommand{\R}{\mathbb{R}}
\newcommand{\E}{\mathbb{E}}
\newcommand{\Prob}{\mathbb{P}}
\newcommand{\T}{^{\top}}
\newcommand{\norm}[1]{\left\lVert #1 \right\rVert}
\newcommand{\normn}[1]{\left\lVert #1 \right\rVert_{n}}
\newcommand{\normL}[1]{\left\lVert #1 \right\rVert_{L_2(P)}}
\newcommand{\Besov}[3]{B^{#1}_{#2,#3}}
\newcommand{\bX}{\mathbf{X}}
\newcommand{\bU}{\mathbf{U}}
\newcommand{\bZ}{\mathbf{Z}}
\newcommand{\btheta}{\boldsymbol{\theta}}
\newcommand{\thetastar}{\btheta^{*}}
\newcommand{\bc}{\mathbf{c}}
\newcommand{\PiJ}{\Pi_{J}}
\newcommand{\VJ}{V_{J}}
\newcommand{\NJ}{N_{J}}
\newtheorem{proposition}{Proposition}
\newtheorem{lemma}{Lemma}
\newtheorem{theorem}{Theorem}
\newtheorem{corollary}{Corollary}
\theoremstyle{definition}
\newtheorem{assumption}{Assumption}
\newtheorem{remark}{Remark}
.
\usepackage{amsmath,amssymb,amsthm}
\newcommand{\R}{\mathbb{R}}
\newcommand{\E}{\mathbb{E}}
\newcommand{\Prob}{\mathbb{P}}
\newcommand{\T}{^{\top}}
\newcommand{\norm}[1]{\left\lVert #1 \right\rVert}
\newcommand{\normn}[1]{\left\lVert #1 \right\rVert_{n}}
\newcommand{\normL}[1]{\left\lVert #1 \right\rVert_{L_2(P)}}
\newcommand{\Besov}[3]{B^{#1}_{#2,#3}}
\newcommand{\bX}{\mathbf{X}}
\newcommand{\bU}{\mathbf{U}}
\newcommand{\bZ}{\mathbf{Z}}
\newcommand{\btheta}{\boldsymbol{\theta}}
\newcommand{\thetastar}{\btheta^{*}}
\newcommand{\bc}{\mathbf{c}}
\newcommand{\PiJ}{\Pi_{J}}
\newcommand{\VJ}{V_{J}}
\newcommand{\NJ}{N_{J}}
\newtheorem{proposition}{Proposition}
\newtheorem{lemma}{Lemma}
\newtheorem{theorem}{Theorem}
\newtheorem{corollary}{Corollary}
\theoremstyle{definition}
\newtheorem{assumption}{Assumption}
\newtheorem{remark}{Remark}
.
\usepackage{amsmath,amssymb,amsthm}
\newcommand{\R}{\mathbb{R}}
\newcommand{\E}{\mathbb{E}}
\newcommand{\Prob}{\mathbb{P}}
\newcommand{\T}{^{\top}}
\newcommand{\norm}[1]{\left\lVert #1 \right\rVert}
\newcommand{\normn}[1]{\left\lVert #1 \right\rVert_{n}}
\newcommand{\normL}[1]{\left\lVert #1 \right\rVert_{L_2(P)}}
\newcommand{\Besov}[3]{B^{#1}_{#2,#3}}
\newcommand{\bX}{\mathbf{X}}
\newcommand{\bU}{\mathbf{U}}
\newcommand{\bZ}{\mathbf{Z}}
\newcommand{\btheta}{\boldsymbol{\theta}}
\newcommand{\thetastar}{\btheta^{*}}
\newcommand{\bc}{\mathbf{c}}
\newcommand{\PiJ}{\Pi_{J}}
\newcommand{\VJ}{V_{J}}
\newcommand{\NJ}{N_{J}}
\newtheorem{proposition}{Proposition}
\newtheorem{lemma}{Lemma}
\newtheorem{theorem}{Theorem}
\newtheorem{corollary}{Corollary}
\theoremstyle{definition}
\newtheorem{assumption}{Assumption}
\newtheorem{remark}{Remark}


\title{E1.3. Aproximação em Besov pelo sieve de wavelets periodizadas}
\author{wafc-draft, \texttt{derivations/02-aproximacao-besov.tex}}
\date{2026-09-18}

\begin{document}
\maketitle

\section*{Onde este resultado entra}

Este arquivo fornece o termo de viés que a desigualdade oráculo (E1.5) e as
taxas (E1.6) consomem: o erro cometido ao substituir cada componente
$\gcomp{\cv}{\md}$ pela sua projeção no espaço de detalhes $\WJ$ da base
periodizada, e o erro correspondente na função de regressão
$f(x,u) = \sum_{\cv} x_{\cv}\,\fcoef{\cv}(u)$. É a transposição do Passo~1
do WALL teórico, com três diferenças: a regularidade é Besov geral
$\Besov{s}{\pi}{r}$, e não $\Besov{s}{2}{\infty}$, para que a regularidade
efetiva $\seff = s - (1/\pi - 1/2)_{+}$ apareça (é ela que separa o sieve
linear do LASSO quando $\pi < 2$); a função de regressão tem o produto
$X_{\cv}\,\gcomp{\cv}{\md}(U_{\md})$, o que troca o fator $p^{2}$ do WALL por
$(pq)^{2}\,C_{X}^{2}$; e a diferença por periodização, que o WALL não trata,
fica registrada na Proposição~\ref{prop:periodizacao}.

Numeração dos resultados: sequência global dos \texttt{derivations/}
(\texttt{docs/TAREFA.md}, §5); este arquivo usa Lema~1, Lema~2, Corolário~1 e
Proposição~1, a acertar pelo chat principal com o que E1.2 tiver numerado.

\section{Objetos}

\paragraph{Base.} Sejam $\phi$ e $\psi$ a função de escala e a wavelet de uma
análise de multirresolução ortonormal de $L_{2}(\R)$, de suporte compacto,
com $\psi$ limitada e com $N$ momentos nulos, $\int y^{a}\psi(y)\,dy = 0$
para $a = 0, \dots, N-1$. Denotamos por $L$ o comprimento do filtro
(\texttt{filter.size} no \texttt{WaveBased}), de modo que
$\operatorname{supp}\psi \subseteq [0, L-1]$ e, para Daublets e Symmlets,
$N = L/2$. A base periodizada em $[0,1]$ é
\[
  \wav{\lev}{\tsl}(u) \;=\; \sum_{\iota \in \mathbb{Z}} 2^{\lev/2}\,
  \psi\bigl(2^{\lev}(u + \iota) - \tsl\bigr),
  \qquad \lev \ge 0,\quad \tsl = 0, \dots, 2^{\lev} - 1,
\]
e o sistema $\{\mathbf{1}\} \cup \{\wav{\lev}{\tsl}\}$ é uma base ortonormal
de $L_{2}[0,1]$ com a medida de Lebesgue, com $\int_{0}^{1}\wav{\lev}{\tsl} = 0$
para toda wavelet (Daubechies, 1992, \S 9.3, ``Wavelets for $L^{p}([0,1])$'',
p.~304). É a base que o \texttt{wbasis()} devolve com
\texttt{boundary = "periodic"} e $\jzero = 0$: a primeira coluna é
$\phi_{00} \equiv 1$, descartada (D2), e as demais são as $\wav{\lev}{\tsl}$,
$\lev = 0, \dots, J-1$.

\paragraph{Sieve e projeção.} $\WJ = \operatorname{span}\{\wav{\lev}{\tsl} :
0 \le \lev < J\}$, de dimensão $\NJ = 2^{J} - 1$, é o complemento ortogonal
das constantes dentro de $\VJ$. Para $g \in L_{2}[0,1]$ com $\int_{0}^{1} g = 0$,
\[
  \PiJ g \;=\; \sum_{\lev = 0}^{J-1}\sum_{\tsl = 0}^{2^{\lev}-1}
  \theta_{\lev\tsl}\,\wav{\lev}{\tsl},
  \qquad \theta_{\lev\tsl} = \inner{g}{\wav{\lev}{\tsl}}
  = \int_{0}^{1} g(u)\,\wav{\lev}{\tsl}(u)\,du,
\]
é a projeção ortogonal de $g$ em $\WJ$, que coincide com a projeção em $\VJ$
porque $\inner{g}{\mathbf{1}} = 0$. Escrevemos $\theta_{\lev\cdot} =
(\theta_{\lev\tsl})_{\tsl}\in\R^{2^{\lev}}$ para o vetor de coeficientes do
nível $\lev$; quando a componente é $\gcomp{\cv}{\md}$, os coeficientes são
os $\coef{\cv}{\md}{\lev}{\tsl}$ da notação.

\paragraph{Regularidade em forma de sequência.} Para $s > 0$ e
$1 \le \pi, r \le \infty$, a norma de Besov em forma de sequência dos
coeficientes de $g$ na base periodizada é
\[
  \norm{g}_{b^{s}_{\pi,r}} \;=\;
  \Bigl\{\sum_{\lev \ge 0}\bigl(2^{\lev(s + 1/2 - 1/\pi)}\,
  \norm{\theta_{\lev\cdot}}_{\pi}\bigr)^{r}\Bigr\}^{1/r},
\]
com a modificação usual para $r = \infty$ (supremo em $\lev$) e para
$\pi = \infty$ (máximo em $\tsl$). É a norma de Donoho e Johnstone (1998,
\S 2) e de Härdle et al.\ (1998, cap.~9, ``Wavelets and Besov Spaces'',
pp.~101--124). A
regularidade efetiva em $L_{2}$ é $\seff = s - (1/\pi - 1/2)_{+}$.

\section{Hipóteses}

\begin{assumption}[base]\label{ass:base}
$\phi$ e $\psi$ são como acima: suporte compacto, $\psi$ limitada, $N$
momentos nulos, e a base periodizada $\{\mathbf{1}\}\cup\{\wav{\lev}{\tsl}\}$
é ortonormal e completa em $L_{2}[0,1]$.
\end{assumption}

\begin{assumption}[regularidade de Besov, forma de sequência]\label{ass:besov}
Existem $s > 0$, $1 \le \pi, r \le \infty$ e $C_{g} < \infty$ tais que, para
todo par $(\cv, \md)$, $\int_{0}^{1}\gcomp{\cv}{\md} = 0$ e
$\norm{\gcomp{\cv}{\md}}_{b^{s}_{\pi,\infty}} \le C_{g}$, isto é,
\[
  \norm{\theta_{\cv\md,\lev\cdot}}_{\pi} \;\le\; C_{g}\,
  2^{-\lev(s + 1/2 - 1/\pi)} \qquad \text{para todo } \lev \ge 0 .
\]
Exige-se $\seff > 0$, ou seja, $s > (1/\pi - 1/2)_{+}$; para o Lema~\ref{lem:sup}
exige-se $s > 1/\pi$.
\end{assumption}

Três observações sobre a Hipótese~\ref{ass:besov}. (i) Ela é tomada como
hipótese primitiva sobre a sequência de coeficientes, como fez o WALL
(Assumption ``Besov regularity''), e não como pertencimento a um espaço de
funções: é a única propriedade de $\gcomp{\cv}{\md}$ que as provas usam, e
evita invocar a equivalência de normas que a base periodizada só fornece
para funções periódicas. (ii) Como $\ell_{r} \subseteq \ell_{\infty}$ com
$\norm{a}_{\infty} \le \norm{a}_{r}$, a hipótese com índice fino $r = \infty$
é a mais fraca da escala, e cobre $\norm{g}_{b^{s}_{\pi,r}} \le C_{g}$ para
qualquer $r$. (iii) Quais funções a satisfazem: se $\psi$ tem $N > s$
momentos nulos e regularidade Hölder maior que $s$, a norma de sequência na
base periodizada é equivalente à norma do espaço de Besov \emph{periódico}
$\Besov{s}{\pi}{r}(\mathbb{T})$, o das funções de $[0,1]$ cuja extensão
$1$-periódica está em $\Besov{s}{\pi}{r}$ localmente (Meyer, 1992, cap.~6, ``Wavelets and spaces of
functions and distributions''; Cohen, 2003, \S 3.6, ``Characterization of
smoothness classes''). A implicação de que a prova
precisa, função em Besov implica decaimento dos coeficientes, usa só os $N > s$
momentos nulos e a limitação de $\psi$; a regularidade de $\psi$ entra só na
recíproca. Uma função de $\Besov{s}{\pi}{r}[0,1]$ que não emenda em $0 \equiv 1$
\emph{não} satisfaz a hipótese para $s > 1/2$; esse caso é a
Proposição~\ref{prop:periodizacao}.

\begin{assumption}[densidade das moduladoras]\label{ass:dens}
$\bU \in [0,1]^{q}$ tem densidade conjunta limitada por $C_{U} < \infty$ (D11).
Em consequência cada marginal $p_{U_{\md}}$ é limitada por $C_{U}$, pois
integrar a conjunta nas outras $q-1$ coordenadas de $[0,1]^{q-1}$ não aumenta
a cota.
\end{assumption}

\begin{assumption}[covariáveis lineares limitadas]\label{ass:X}
$\max_{\cv}|X_{\cv}| \le C_{X}$ quase certamente.
\end{assumption}

\section{O lema em $L_{2}$}

\begin{lemma}[erro de projeção em $L_{2}$]\label{lem:L2}
Sob as Hipóteses~\ref{ass:base} e~\ref{ass:besov}, para todo $(\cv,\md)$ e
todo $J \ge 0$,
\[
  \norm{\gcomp{\cv}{\md} - \PiJ\gcomp{\cv}{\md}}_{L_{2}[0,1]}^{2}
  \;=\; \sum_{\lev \ge J}\norm{\theta_{\cv\md,\lev\cdot}}_{2}^{2}
  \;\le\; \frac{C_{g}^{2}}{1 - 2^{-2\seff}}\; 2^{-2J\seff}.
\]
Se além disso vale a Hipótese~\ref{ass:dens},
\[
  \norm{\gcomp{\cv}{\md} - \PiJ\gcomp{\cv}{\md}}_{L_{2}(P_{U_{\md}})}^{2}
  \;=\; \E\bigl[(\gcomp{\cv}{\md} - \PiJ\gcomp{\cv}{\md})^{2}(U_{\md})\bigr]
  \;\le\; \frac{C_{U}\,C_{g}^{2}}{1 - 2^{-2\seff}}\; 2^{-2J\seff}.
\]
\end{lemma}

\begin{proof}
Fixe $(\cv,\md)$ e escreva $g = \gcomp{\cv}{\md}$, $\theta_{\lev\tsl} =
\coef{\cv}{\md}{\lev}{\tsl}$.

\emph{Passo 1 (soma de cauda).} A Hipótese~\ref{ass:besov} dá
$\norm{\theta_{\lev\cdot}}_{\pi} < \infty$ em cada nível; pelo Passo~2 abaixo
$\sum_{\lev}\norm{\theta_{\lev\cdot}}_{2}^{2} < \infty$, logo $g \in
L_{2}[0,1]$ e, como a base é completa e $\inner{g}{\mathbf{1}} = 0$,
$g = \sum_{\lev\ge 0}\sum_{\tsl}\theta_{\lev\tsl}\wav{\lev}{\tsl}$ em $L_{2}$.
A projeção ortogonal em $\WJ$ retém os níveis $\lev < J$, e Parseval dá
\[
  \norm{g - \PiJ g}_{L_{2}[0,1]}^{2}
  = \sum_{\lev\ge J}\sum_{\tsl=0}^{2^{\lev}-1}\theta_{\lev\tsl}^{2}
  = \sum_{\lev \ge J}\norm{\theta_{\lev\cdot}}_{2}^{2}.
\]

\emph{Passo 2 (de $\ell_{\pi}$ para $\ell_{2}$ dentro do nível).} Em
$\R^{2^{\lev}}$: se $\pi \le 2$, $\norm{a}_{2} \le \norm{a}_{\pi}$; se
$\pi > 2$, a desigualdade de Hölder dá $\norm{a}_{2} \le
(2^{\lev})^{1/2 - 1/\pi}\norm{a}_{\pi}$. Nos dois casos
$\norm{\theta_{\lev\cdot}}_{2} \le 2^{\lev(1/2 - 1/\pi)_{+}}
\norm{\theta_{\lev\cdot}}_{\pi}$, e com a Hipótese~\ref{ass:besov}
\[
  \norm{\theta_{\lev\cdot}}_{2}
  \;\le\; C_{g}\,2^{-\lev\{s + 1/2 - 1/\pi - (1/2 - 1/\pi)_{+}\}}
  \;=\; C_{g}\,2^{-\lev\seff},
\]
pois $s + 1/2 - 1/\pi - (1/2 - 1/\pi)_{+}$ vale $s$ quando $\pi \ge 2$ e
$s - (1/\pi - 1/2)$ quando $\pi < 2$. É aqui que a perda $(1/\pi - 1/2)_{+}$
aparece: para $\pi < 2$ a norma $\ell_{\pi}$ do nível controla a $\ell_{2}$
sem ganho, e o expoente de $L_{2}$ fica $\seff < s$.

\emph{Passo 3 (soma geométrica).}
$\sum_{\lev\ge J}C_{g}^{2}\,2^{-2\lev\seff} = C_{g}^{2}\,2^{-2J\seff}/(1 - 2^{-2\seff})$,
finita porque $\seff > 0$.

\emph{Passo 4 (densidade).} $\E[h^{2}(U_{\md})] = \int_{0}^{1}h^{2}\,p_{U_{\md}}
\le C_{U}\norm{h}_{L_{2}[0,1]}^{2}$ com $h = g - \PiJ g$.
\end{proof}

\section{O lema em $L_{\infty}$}

O Lema~\ref{lem:L2} basta para o termo de viés da desigualdade oráculo. A
versão uniforme é o que o WALL usa na localização (Assumption ``Besov
regularity'', segunda consequência) e o que E1.4 e E1.5 podem precisar para
controlar $\norm{f - f_{J}}_{n}$ de forma determinística ou para limitar as
colunas do desenho.

\begin{lemma}[erro de projeção uniforme]\label{lem:sup}
Sob as Hipóteses~\ref{ass:base} e~\ref{ass:besov} com $s > 1/\pi$, cada
$\gcomp{\cv}{\md}$ tem versão contínua e, para todo $J \ge 0$,
\[
  \norm{\gcomp{\cv}{\md} - \PiJ\gcomp{\cv}{\md}}_{\infty}
  \;\le\; \frac{(L-1)\,\norm{\psi}_{\infty}\,C_{g}}{1 - 2^{-(s - 1/\pi)}}\;
  2^{-J(s - 1/\pi)} .
\]
\end{lemma}

\begin{proof}
\emph{Sobreposição limitada.} Para $u \in [0,1]$ e $\lev \ge 0$,
\[
  \sum_{\tsl=0}^{2^{\lev}-1}|\wav{\lev}{\tsl}(u)|
  \;\le\; \sum_{\tsl\in\mathbb{Z}} 2^{\lev/2}\,|\psi(2^{\lev}u - \tsl)|
  \;\le\; (L-1)\,2^{\lev/2}\,\norm{\psi}_{\infty},
\]
porque a soma sobre as translações periodizadas é a soma sobre todas as
translações inteiras em $\R$, e $\psi(2^{\lev}u - \tsl) \ne 0$ exige
$2^{\lev}u - \tsl \in (0, L-1)$, o que no máximo $L - 1$ inteiros $\tsl$
satisfazem.

\emph{Cota por nível.} Com $\max_{\tsl}|\theta_{\lev\tsl}| \le
\norm{\theta_{\lev\cdot}}_{\pi} \le C_{g}\,2^{-\lev(s + 1/2 - 1/\pi)}$,
\[
  \Bigl\lVert \sum_{\tsl}\theta_{\lev\tsl}\wav{\lev}{\tsl}\Bigr\rVert_{\infty}
  \;\le\; \max_{\tsl}|\theta_{\lev\tsl}| \sup_{u}\sum_{\tsl}|\wav{\lev}{\tsl}(u)|
  \;\le\; (L-1)\norm{\psi}_{\infty}\,C_{g}\,2^{-\lev(s - 1/\pi)} .
\]

\emph{Soma.} Como $s > 1/\pi$, a série $\sum_{\lev}\sum_{\tsl}\theta_{\lev\tsl}
\wav{\lev}{\tsl}$ converge absoluta e uniformemente; o limite uniforme é
contínuo e coincide q.c.\ com o limite em $L_{2}$, que é $g$ (Passo~1 do
Lema~\ref{lem:L2}). Logo $g$ tem versão contínua, $g - \PiJ g$ é a soma
uniforme dos níveis $\lev \ge J$, e a soma geométrica das cotas por nível dá o
enunciado.
\end{proof}

Para $\pi = 2$ a taxa é $2^{-J(s - 1/2)}$, a do WALL. A perda de $L_{2}$ para
$L_{\infty}$ é $1/\pi$ e não $(1/\pi - 1/2)_{+}$: o controle uniforme paga
o encaixe $\Besov{s}{\pi}{r} \hookrightarrow \Besov{s - 1/\pi}{\infty}{r}$
inteiro.

\section{O viés da função de regressão}

Seja $f(x,u) = \sum_{\cv=1}^{p} x_{\cv}\bigl\{\lvl{\cv} +
\sum_{\md=1}^{q}\gcomp{\cv}{\md}(u_{\md})\bigr\}$ a função de regressão e
\[
  f_{J}(x,u) \;=\; \sum_{\cv=1}^{p} x_{\cv}\Bigl\{\lvl{\cv} +
  \sum_{\md=1}^{q}(\PiJ\gcomp{\cv}{\md})(u_{\md})\Bigr\}
  \;=\; \sum_{\cv} x_{\cv}\lvl{\cv} + \sum_{\cv,\md}\sum_{\lev<J}\sum_{\tsl}
  \coefs{\cv}{\md}{\lev}{\tsl}\; x_{\cv}\,\wav{\lev}{\tsl}(u_{\md})
\]
a sua aproximação no sieve, com $\coefs{\cv}{\md}{\lev}{\tsl} =
\inner{\gcomp{\cv}{\md}}{\wav{\lev}{\tsl}}$. O vetor $(\bc, \thetastar)$ é
o oráculo de E1.5: $f_{J}(\bX_{i},\bU_{i}) = \bZ_{i}\T(\bc,\thetastar)$
com a matriz de desenho da notação (ordem das colunas D12).

\begin{corollary}[viés do sieve]\label{cor:vies}
Sob as Hipóteses~\ref{ass:base} a~\ref{ass:X},
\[
  \norm{f - f_{J}}_{L_{2}(P)}^{2}
  \;=\; \E\bigl[(f - f_{J})^{2}(\bX,\bU)\bigr]
  \;\le\; \frac{(pq)^{2}\,C_{X}^{2}\,C_{U}\,C_{g}^{2}}{1 - 2^{-2\seff}}\;
  2^{-2J\seff},
\]
e, se $s > 1/\pi$, para todo $(x,u) \in [-C_{X},C_{X}]^{p}\times[0,1]^{q}$,
\[
  |f(x,u) - f_{J}(x,u)| \;\le\;
  \frac{pq\,C_{X}\,(L-1)\norm{\psi}_{\infty}\,C_{g}}{1 - 2^{-(s-1/\pi)}}\;
  2^{-J(s - 1/\pi)} ,
\]
de modo que $\norm{f - f_{J}}_{n} \le \norm{f - f_{J}}_{\infty}$ obedece à
mesma cota, para qualquer amostra.
\end{corollary}

\begin{proof}
$f - f_{J} = \sum_{\cv,\md} X_{\cv}\,h_{\cv\md}(U_{\md})$ com $h_{\cv\md} =
\gcomp{\cv}{\md} - \PiJ\gcomp{\cv}{\md}$; os níveis $\lvl{\cv}$ cancelam.
Cauchy--Schwarz sobre os $pq$ termos e $|X_{\cv}| \le C_{X}$ dão
\[
  \E\bigl[(f - f_{J})^{2}\bigr]
  \;\le\; pq\sum_{\cv,\md}\E\bigl[X_{\cv}^{2}\,h_{\cv\md}^{2}(U_{\md})\bigr]
  \;\le\; pq\,C_{X}^{2}\sum_{\cv,\md}\E\bigl[h_{\cv\md}^{2}(U_{\md})\bigr],
\]
e o Lema~\ref{lem:L2} em cada um dos $pq$ termos fecha a primeira cota: um
fator $pq$ vem de Cauchy--Schwarz e outro da soma, como o $p^{2}$ do WALL. A
segunda é a desigualdade triangular com o Lema~\ref{lem:sup}.
\end{proof}

A dependência entre $\bX$ e $\bU$ não entra: o viés só usa a limitação de
$\bX$ e a densidade marginal de cada $U_{\md}$. O que a dependência afeta é
a passagem da norma de predição para a norma de cada componente, que é a
condição de desenho de E1.4.

\section{A diferença por periodização}

O WALL adota a hipótese em forma de sequência e não discute o que ela exclui.
Para o WAFC vale registrar, porque a aplicação terá moduladoras cujos efeitos
não emendam nas bordas (renda, idade, temperatura), e porque é o que motiva
o reescalonamento para $[\varepsilon, 1-\varepsilon]$ do \texttt{wall()} e a
opção \texttt{boundary = "interval"}.

\begin{proposition}[custo da periodização]\label{prop:periodizacao}
Seja a base da Hipótese~\ref{ass:base} e seja $g \in \Besov{s}{\pi}{r}[0,1]$
(espaço de Besov do intervalo, sem condição de periodicidade) com
$s > 1/\pi$, $N > s$ e $\int_{0}^{1} g = 0$. Então existe $C$ que depende de
$g$, de $\psi$ e de $(s,\pi)$ tal que, para todo $J$ com $2^{J} \ge L - 1$,
\[
  \norm{g - \PiJ g}_{L_{2}[0,1]}^{2}
  \;\le\; C\,2^{-2J\seff} \;+\; 2(L-1)\,\norm{\psi}_{1}^{2}\,\norm{g}_{\infty}^{2}\; 2^{-J},
\]
isto é, a taxa em $L_{2}$ é $2^{-J\min(\seff,\,1/2)}$. Quando $g(0) \ne g(1)$
o segundo termo é da ordem exata $2^{-J}$ e $\norm{g - \PiJ g}_{\infty}$
não converge a zero.
\end{proposition}

\begin{proof}[Prova (esboço; a extensão é o passo a verificar)]
Seja $\tilde g$ a extensão $1$-periódica de $g$ a $\R$. Por construção da
base periodizada, $\theta_{\lev\tsl} = \int_{\R}\tilde g(y)\,2^{\lev/2}
\psi(2^{\lev}y - \tsl)\,dy$. Para $2^{\lev} \ge L-1$, separe em cada nível as
translações \emph{interiores}, aquelas com $[\tsl, \tsl + L - 1] \subseteq
[0, 2^{\lev}]$, cujo suporte fica dentro de $[0,1]$, das translações de
\emph{borda}, que atravessam $0 \equiv 1$; há no máximo $L - 1$ destas por
nível.

\emph{Interiores.} Nelas $\theta_{\lev\tsl} = \inner{Eg}{\psi_{\lev\tsl}}_{\R}$
para qualquer extensão $Eg$ de $g$ a $\R$. Tomando $E$ o operador de
extensão limitado de $\Besov{s}{\pi}{r}[0,1]$ em $\Besov{s}{\pi}{r}(\R)$
(Triebel, 1983, cap.~3, ``Function Spaces on Domains'', pp.~188--211;
Cohen, 2003, \S 3.9, ``Bounded domains'') e usando a
caracterização de $\Besov{s}{\pi}{r}(\R)$ pelos coeficientes de wavelet, que
para a implicação função $\Rightarrow$ sequência exige só $N > s$ momentos
nulos, obtém-se $\norm{(\theta_{\lev\tsl})_{\tsl\ \mathrm{int}}}_{\pi} \le
C_{1}\,2^{-\lev(s + 1/2 - 1/\pi)}$, e a soma de cauda dos interiores é
$\le C\,2^{-2J\seff}$ pelo argumento do Lema~\ref{lem:L2}.

\emph{Borda.} Como $s > 1/\pi$, $g$ é limitada, e $|\theta_{\lev\tsl}| \le
\norm{g}_{\infty}\norm{\psi_{\lev\tsl}}_{1} = \norm{g}_{\infty}\norm{\psi}_{1}
2^{-\lev/2}$; com $L - 1$ termos por nível,
$\sum_{\lev\ge J}\sum_{\tsl\ \mathrm{borda}}\theta_{\lev\tsl}^{2} \le
(L-1)\norm{\psi}_{1}^{2}\norm{g}_{\infty}^{2}\sum_{\lev\ge J}2^{-\lev} =
2(L-1)\norm{\psi}_{1}^{2}\norm{g}_{\infty}^{2}\,2^{-J}$.

\emph{Ordem exata.} Se $g$ é contínua em $[0,1]$ e $g(0) \ne g(1)$, para
uma translação de borda com $2^{\lev}$ grande $\tilde g$ é, no suporte de
$\psi_{\lev\tsl}$, $g(1) + o(1)$ à esquerda de $0$ e $g(0) + o(1)$ à direita,
logo $\theta_{\lev\tsl} = 2^{-\lev/2}\{g(0) - g(1)\}\int_{y > a_{\tsl}}\psi(y)\,dy
+ o(2^{-\lev/2})$, com $a_{\tsl}$ o ponto do suporte de $\psi$ que cai em
$0$; como $\int\psi = 0$ e $\psi$ não é nula em nenhum intervalo do suporte,
essa integral parcial é não nula para algum $\tsl$ e o termo de borda é de
ordem $2^{-\lev}$ em cada nível. A não convergência uniforme é a de qualquer
série de wavelets num salto: $\PiJ g$ é contínua e $\tilde g$ não.
\end{proof}

O que isso muda para o WAFC:
\begin{itemize}[nosep]
  \item Sob a Hipótese~\ref{ass:besov} nada muda: ela já é uma hipótese
    periódica. O que a Proposição diz é o \emph{preço} de a hipótese falhar:
    com uma $\gcomp{\cv}{\md}$ que não emenda, o viés cai como $2^{-J/2}$, o
    que em E1.6 leva a taxa de predição a $n^{-1/2}$ (a menos de logaritmos)
    em vez de $n^{-2\seff/(2\seff+1)}$, uma perda para todo $\seff > 1/2$.
  \item Com \texttt{boundary = "interval"} (Cohen, Daubechies e Vial, 1993),
    o espaço $\VJ$ do intervalo reproduz polinômios de grau $< N$ e a sua
    base caracteriza $\Besov{s}{\pi}{r}[0,1]$ sem periodicidade (Cohen,
    2003, \S 3.9); o Lema~\ref{lem:L2} vale então com
    $\Besov{s}{\pi}{r}[0,1]$ no lugar do espaço periódico e a mesma prova,
    pois $\mathbf{1} \in V_{\jzero}$ e a projeção de uma função de média
    zero em $\VJ$ é ortogonal a $\mathbf{1}$. O preço é numérico e
    estrutural: o bloco de escala $V_{\jzero}$ tem $2^{\jzero}$ funções, com
    $\jzero$ grande o bastante para as duas zonas de borda não se tocarem
    (o \texttt{wbasis()} exige $\jzero \ge 4$ para $L = 8$), a constante
    deixa de ser uma coluna isolada e passa a ser uma combinação das funções
    de escala, e a restrição $\int\gcomp{\cv}{\md} = 0$ tem de ser imposta de
    outro modo (E1.2, E2.1).
  \item O reescalonamento para $[\varepsilon, 1-\varepsilon]$ do
    \texttt{wall()} ($\varepsilon = 1.9^{-J}$; Montoril, Pinheiro e
    Vidakovic, 2019) afasta os dados da zona onde a periodização deixa
    artefato, mas não entra na teoria (D11). Note que o termo de viés que
    E1.5 usa é $\inf$ sobre o sieve da norma $L_{2}(P_{U})$, e não
    necessariamente $\PiJ g$: quando $P_{U}$ não visita a zona de borda, o
    melhor elemento de $\WJ$ em $L_{2}(P_{U})$ pode ser bem melhor que a
    projeção de Lebesgue. Isso fica como pergunta (handoff).
\end{itemize}

\section{O que foi transposto e o que é novo}

\begin{itemize}[nosep]
  \item Transposto do WALL (Passo~1 e prova S.1): a hipótese em forma de
    sequência; a soma de cauda por Parseval; a soma geométrica; o
    Cauchy--Schwarz sobre as componentes com a densidade limitada.
  \item Generalizado: $\Besov{s}{\pi}{r}$ com $\pi$ qualquer, via o Passo~2
    do Lema~\ref{lem:L2}, de onde sai $\seff$; a versão uniforme com a
    constante explícita de sobreposição $(L-1)\norm{\psi}_{\infty}$.
  \item Novo: o produto $X_{\cv}\,\gcomp{\cv}{\md}(U_{\md})$ no
    Corolário~\ref{cor:vies}, que custa $C_{X}^{2}$ e troca $p^{2}$ por
    $(pq)^{2}$; a Proposição~\ref{prop:periodizacao}.
\end{itemize}

\section{O que isto não cobre}

\begin{itemize}[nosep]
  \item A aproximação \emph{não linear}. Para $\pi < 2$ a melhor aproximação
    com $M$ termos escolhidos livremente atinge $M^{-s}$, enquanto o sieve
    linear atinge só $\NJ^{-\seff}$; é essa diferença que o LASSO explora e
    que o corolário de compressibilidade (E1.6, weak-$\ell_{\tau}$ com
    $\tau = (s + 1/2)^{-1}$ no caso crítico) formaliza. O lema aqui é só o
    viés do sieve cheio.
  \item A norma empírica $\normn{f - f_{J}}$ além da cota uniforme: a versão
    com probabilidade alta a partir de $\norm{f - f_{J}}_{L_{2}(P)}$ é
    concentração e fica em E1.5.
  \item A regularidade diferente por bloco ($s_{\cv\md}$, $\pi_{\cv\md}$): o
    enunciado tomou $(s, \pi, C_{g})$ comuns; a versão heterogênea é a mesma
    prova com $\min_{\cv\md}\seff_{\cv\md}$ na taxa.
  \item A recíproca da caracterização (sequência $\Rightarrow$ função), que
    exigiria regularidade Hölder de $\psi$ maior que $s$, e que as provas de
    E1 não usam.
  \item A extensão de $\Besov{s}{\pi}{r}[0,1]$ a $\R$ na
    Proposição~\ref{prop:periodizacao}, que está citada e não provada.
\end{itemize}

\section{Conferência numérica}

\texttt{check/02-aproximacao-besov.R} (imprime \texttt{OK}; cerca de 50~s).
Base Daublets com $L = 8$ ($N = 4$), $\jzero = 0$, constante descartada;
projeção por mínimos quadrados numa grade de $2^{12}$ pontos ($J \le 9$) e
coeficientes por quadratura numa grade de $2^{15}$ pontos ($\lev \le 12$). O
que saiu:
\begin{itemize}[nosep]
  \item Parseval: o erro de projeção em $\WJ$ e a soma de cauda dos
    coeficientes diferem em no máximo $1.4\times10^{-4}$ relativo (bumps,
    $J = 1, \dots, 8$).
  \item $\sin(2\pi u)$, periódica e $C^{\infty}$: a queda por nível é
    $3.92, 3.98, 4.00, 4.00, 4.00$ bits em $L_{2}$ e o mesmo em $L_{\infty}$
    ($J = 4, \dots, 8$), isto é, $2^{-JN}$: para uma função periódica lisa o
    limite é o número de momentos nulos, como a prova prevê.
  \item bumps (Donoho e Johnstone), com $\seff = 3/2$ pelas quinas
    $|u - t_{i}|$: as normas por nível $\norm{\theta_{\lev\cdot}}_{2}$ caem
    $1.49, 1.35, 1.27$ bits em $\lev = 9 \to 12$ (média $1.37$; tolerância
    $1/4$). Abaixo de $\lev \approx 8$ a queda é menor que $1$ bit porque a
    largura mínima das protuberâncias é $0.005 \approx 2^{-7.6}$; é o regime
    pré-assintótico, e é também a situação em que o LASSO ganha do sieve.
    Uma quina periódica isolada, $|\sin\pi(u - 0.4)|$, dá quedas que oscilam
    com a fase da quina na grade diádica ($2.66, 1.56, 0.13, 1.65, \dots$) e
    média de quatro níveis igual a $1.50$.
  \item Funções que não emendam ($u - 1/2$ e $e^{u}$ centrada) na base
    periodizada: queda em $L_{2}$ de $0.50, 0.50, 0.51$ bits por nível em
    $J = 7, 8, 9$, independentemente da regularidade interior, e erro
    $L_{\infty}$ estacionado em $0.54$ e $0.92$ (saltos de $1.00$ e $1.72$):
    é a Proposição~\ref{prop:periodizacao}.
  \item As mesmas funções na base do intervalo ($\jzero = 4$, $\VJ$
    inteiro, $J = 4, \dots, 9$): a linear é reproduzida a $10^{-10}$ e a
    exponencial cai de $6.3\times10^{-6}$ em $J = 4$ a $9.3\times10^{-11}$ em
    $J = 9$ (cerca de $4.4$ bits por nível até o piso numérico), contra
    $0.13$ e $0.023$ na periodizada.
\end{itemize}

\section*{Referências}
\sloppy

Todas as referências deste arquivo foram conferidas em L3 (2026-09-19) e
estão em \texttt{docs/referencias-verificadas.bib}. A única localização que a
rodada não conseguiu confirmar está anotada abaixo, em Meyer (1992).

\begin{itemize}[nosep]
  \item Cohen, A. (2003). \emph{Numerical Analysis of Wavelet Methods}.
    Studies in Mathematics and its Applications 32. Elsevier.
    (\texttt{Cohen-2003}.) O cap.~3 é ``Approximation and smoothness''; a
    caracterização das classes de suavidade é a \S 3.6 e os domínios
    limitados, com a extensão, são a \S 3.9.
  \item Cohen, A., Daubechies, I. e Vial, P. (1993). Wavelets on the
    interval and fast wavelet transforms. \emph{Applied and Computational
    Harmonic Analysis} 1, 54--81. (\texttt{cohen1993wavelets})
  \item Daubechies, I. (1992). \emph{Ten Lectures on Wavelets}. CBMS-NSF
    Regional Conference Series in Applied Mathematics 61. SIAM.
    (\texttt{Daubechies-1992}.) As wavelets periodizadas são a \S 9.3,
    ``Wavelets for $L^{p}([0,1])$'', p.~304.
  \item Donoho, D. L. e Johnstone, I. M. (1998). Minimax estimation via
    wavelet shrinkage. \emph{The Annals of Statistics} 26, 879--921.
    (\texttt{Donoho-Johnstone-1998})
  \item Härdle, W., Kerkyacharian, G., Picard, D. e Tsybakov, A. (1998).
    \emph{Wavelets, Approximation, and Statistical Applications}. Lecture
    Notes in Statistics 129. Springer. (\texttt{hardle1998wavelets},
    copiada do WALL.) O cap.~9 é ``Wavelets and Besov Spaces'', pp.~101--124.
  \item Meyer, Y. (1992). \emph{Wavelets and Operators}, volume~1. Cambridge
    Studies in Advanced Mathematics 37. Cambridge University Press.
    (\texttt{Meyer-1992}.) O cap.~6 é ``Wavelets and spaces of functions and
    distributions'', pp.~163--207. A citação original deste arquivo trazia
    ``cap.~III, \S 11'' para a base periodizada; L3 não conseguiu conferir as
    seções internas dos capítulos (sem acesso ao volume) e trocou aquela
    citação por Daubechies (1992, \S 9.3), que foi conferida.
  \item Montoril, M. H., Pinheiro, A. e Vidakovic, B. (2019). Wavelet-based
    estimators for mixture regression. \emph{Scandinavian Journal of
    Statistics} 46, 215--234. (\texttt{Montoril-Pinheiro-Vidakovic-2019})
  \item Triebel, H. (1983). \emph{Theory of Function Spaces}. Monographs in
    Mathematics 78. Birkhäuser, Basel. (\texttt{Triebel-1983}.) O cap.~3 é
    ``Function Spaces on Domains'', pp.~188--211.
\end{itemize}

\end{document}
